% ============================================================================
% Rewritten study text (drop-in snippets). Tables referenced here live in
% ../study2_tables/ and are regenerated by scripts/make_study2_tables.py.
% ============================================================================

% ---------------------------------------------------------------------------
% STUDY 1a — replaces the "130x / identical to numerical precision" framing
% ---------------------------------------------------------------------------
\subsection{Study 1A: permutation invariance (a correctness check) (A1)}
\label{sec:study1a}
\label{sec:A1}

\paragraph{Why it matters.}
A neuron permutation is a quiver isomorphism, so knowledge-matrix
invariance under it is automatic by construction
(Proposition~\ref{prop:perm} is a corollary of quiver-isomorphism
invariance, not a hard-won discovery). This experiment is therefore a
numerical \emph{correctness check}: it confirms that the implementation
realises the identity the theory guarantees, and quantifies the float32
accumulation that separates the computed object from its exact value. It
is not the paper's primary evidence for invariance; the load-bearing,
non-automatic invariance is the cross-architecture case
(Theorem~\ref{thm:maxinv}, Study~3), whose only transform-based evidence
is the approximate teleportation of Study~1B.

\paragraph{Setup.}
Architecture: ResNet-152 (the only trio member that admits an in-place
channel permutation; DenseNet-121 and GoogLeNet have concatenation
topologies with no in-place channel swap and are covered by
teleportation instead). Transform: random permutations of the
\emph{wide-face} permutation subgroup --- the 2048-channel post-pool
faces; bottleneck interiors are not permutable in place. Inputs for the
adversarial signal: clean/adversarial pairs from the six-attack suite
(Section~\ref{sec:setup}). We measure both drifts in raw fp64 Frobenius
(KM) and $\ell_2$ (penultimate), reported \emph{relative to the same
model's adversarial-pair signal} so the comparison is unit-free.

\paragraph{Procedure.}
\begin{enumerate}
\item Compute $\M(x)$ and $h(x)$ for each input.
\item Draw a random wide-face neuron permutation $P$ and build the
isomorphic network $\tilde f = P \cdot f$ (identical function, permuted
weights).
\item Recompute $\M_{\tilde f}(x)$ and $h_{\tilde f}(x)$.
\item Record the permutation \emph{noise} $\|\M_{\tilde f}(x)-\M(x)\|_F$
and $\|h_{\tilde f}(x)-h(x)\|_2$.
\item Record the adversarial \emph{signal} $\|\M(x')-\M(x)\|_F$ and
$\|h(x')-h(x)\|_2$ over the attack pairs, and report each noise relative
to its own signal.
\end{enumerate}

\paragraph{What we measure and the claim it licenses.}
The signal-to-noise ratio per representation. The claim it licenses is
narrow and correct: the knowledge matrix is permutation-invariant up to
float32 accumulation, whereas a statistic of the permuted penultimate
features moves by an amount of the same order as the adversarial signal
it is meant to register --- exactly as Theorem~\ref{thm:maxinv} requires.
It does \emph{not} license a superiority claim: a permutation-equivariant
penultimate statistic would also be invariant, and the gradient$\times$input
family shares the knowledge matrix's invariance (Theorem~\ref{thm:weff}).

\paragraph{Result.}
Measured signal-relative (Table~\ref{tab:signal-relative}), the
knowledge-matrix permutation noise is $0.1$--$0.2\%$ of its
adversarial signal in the mean (the signal sits
$7{\times}10^{2}$--$1.7{\times}10^{7}$ above the noise floor across
architectures and attacks), whereas the penultimate drift is of the
same order as its own adversarial signal (at most $3.05\times$ below it,
and \emph{above} it in 16 of 24 cells). Honesty rows we commit to: the
worst-case tail ratio on ResNet-152 (minimum signal over maximum noise)
reaches $1.14\times$ for DeepFool --- the extreme tails nearly touch on
the BN-heaviest architecture; the residual KM drift (mean $0.116$, max
$5.78$ Frobenius) is float32 accumulation, not a violation of the
theorem --- float64 spot checks reduce it to numerical zero. DenseNet-121
and GoogLeNet admit no in-place channel permutation (concatenation
topologies) and are covered by teleportation instead.

\paragraph{Reading.}
The result shows that the implementation realises the construction's
guaranteed permutation invariance to float32 accumulation, and that a
penultimate statistic does not. It does \emph{not} show that this
invariance is unique to knowledge matrices, nor that the matrix is a
better representation --- only that the identity holds where the theory
says it must, on the one trio member where the permutation is realisable
in place.

% ---------------------------------------------------------------------------
% STUDY 1b — replaces the theorem-asserted zero
% ---------------------------------------------------------------------------
\subsection{Study 1B: teleportation, under the three-tier claim structure (A2)}
\label{sec:study1b}
\label{sec:A2}

\paragraph{Why it matters.}
Teleportation is the only \emph{multi-architecture} invariance arm and the
only one that exercises a transform broader than a quiver isomorphism. But
on these BatchNorm networks it is function-\emph{approximate}, so the
invariance theorem's hypothesis does not hold and no table entry may be
asserted from it. This study is the honest reporting of that situation:
it measures what can be measured and bounds only what the violated
hypothesis still permits.

\paragraph{Setup.}
Architectures: all three trio members (ResNet-152, DenseNet-121,
GoogLeNet). Transform: neural teleportation \citep{armenta2024neural} ---
a function-preserving change-of-basis (per-neuron rescaling) realised by
the \texttt{neuralteleportation} library's COB models. Canonical run:
$T = 50$ random teleportations per architecture (seed indices
$0,\ldots,49$), evaluated on $N = 25{,}000$ ImageNet-validation samples
(the first $25{,}000$ by sorted filename), matching the compiled panel
(Table~\ref{tab:s1_table}, Section~\ref{sec:study1c}). We verify
function-equivalence per teleportation via the $10^{-3}$ logit-drift
gate.

\paragraph{Procedure.}
\begin{enumerate}
\item Compute $\M(x)$, $h(x)$, $f(x)$ for each sample.
\item Draw a random COB teleportation $\tau$ and build $\tilde f = \tau
\cdot f$, an approximate-isomorphism image.
\item Measure the logit gate $\varepsilon = \|f_\tau(x) - f(x)\|$ per
sample and test it against $10^{-3}$.
\item Compute the \emph{visible} knowledge-matrix drift
$\|(\M_\tau(x)-\M(x))\mathbf 1\|$ and verify per pair that it equals the
gate $\varepsilon$ (Proposition~\ref{prop:visdrift}(ii)).
\item Recompute the penultimate features $h_\tau(x)$ and pass the pairs
to the similarity panel of Study~1C (Section~\ref{sec:study1c}).
\end{enumerate}

\paragraph{What we measure and the claim it licenses.}
The logit gate, the visible knowledge-matrix drift (an \emph{identity}
equal to the gate by the row-sum identity, not a measurement of
invariance), and the penultimate drift that the Study~1C panel
adjudicates. The claim this licenses is exactly the three-tier structure
of Proposition~\ref{prop:visdrift}: under an exact isomorphism the drift
would be zero; under this approximate teleportation the visible component
equals the gate; the invisible component carries no a priori bound from
the gate (Theorem~\ref{thm:nogo}) and is \emph{not} estimated here.

\paragraph{Result.}
In our runs, $100/100$ teleportations exceed the $10^{-3}$ logit-drift
gate (up to $10^{-1}$ on ResNet-152, under $4\times10^{-2}$ otherwise),
so the invariance theorem's hypothesis does not hold and no table entry
may be asserted from it. We report the knowledge-matrix drift through the
visible/invisible split: the visible channel obeys the equality
$\|(\M_\tau(x)-\M(x))\mathbf 1\|=\|f_\tau(x)-f(x)\|$, which we verify per
pair against the logit gate; the invisible component is left unestimated.
The penultimate features drift substantially under teleportation (the
panel of Study~1C reports Procrustes shape distance $632$--$1471$ and
soft-matching distance $22$--$61$ across the trio,
Table~\ref{tab:s1_table}); against these the panel adjudicates which
measures quotient the drift out.

\paragraph{Reading.}
This study does \emph{not} verify invariance across the three
architectures: its transform fails the isomorphism hypothesis $100/100$.
What it shows is the visible knowledge-matrix drift pinned to the
function-approximation gate by an identity, with the invisible component
honestly flagged as unbounded and unmeasured. We credit the three-tier
framing as the correct reporting of an approximate transform, not as
evidence the matrix is invariant under one.
